\documentclass[aspectratio=169]{beamer}

% Match the existing CSE 534 Beamer styling used in the N-gram deck
\usetheme{default}
\usecolortheme{default}
\setbeamertemplate{navigation symbols}{}
\setbeamertemplate{footline}[frame number]

\definecolor{darkbg}{HTML}{000000}
\definecolor{lighttext}{HTML}{999999}
\definecolor{miamired}{HTML}{941728}
\definecolor{codebg}{HTML}{1a1a1a}

\setbeamercolor{background canvas}{bg=darkbg}
\setbeamercolor{normal text}{fg=lighttext}
\setbeamercolor{frametitle}{fg=miamired}
\setbeamercolor{title}{fg=miamired}
\setbeamercolor{structure}{fg=miamired}
\setbeamercolor{item}{fg=miamired}
\setbeamercolor{block title}{fg=miamired,bg=codebg}
\setbeamercolor{block body}{fg=lighttext,bg=codebg}

\setbeamerfont{title}{series=\bfseries,size=\Large}
\setbeamerfont{frametitle}{series=\bfseries}

\usepackage{graphicx}
\usepackage{listings}
\usepackage{tikz}
\usepackage{array}
\usepackage{booktabs}
\usepackage{colortbl}
\usepackage{amsmath}
\usepackage{amssymb}

\lstset{
    basicstyle=\ttfamily\small\color{lighttext},
    backgroundcolor=\color{codebg},
    keywordstyle=\color[HTML]{569cd6},
    commentstyle=\color[HTML]{6a9955},
    stringstyle=\color[HTML]{ce9178},
    showstringspaces=false,
    breaklines=true,
    frame=single,
    rulecolor=\color{codebg}
}

\setlength{\arrayrulewidth}{0.5pt}
\arrayrulecolor{miamired}
\renewcommand{\textbf}[1]{\textcolor{white}{\bfseries #1}}

\title{Entropy, Cross-Entropy, and KL Divergence}
\subtitle{Information theory for generative AI}
\author{CSE 534}
\date{}

\begin{document}

\begin{frame}
\titlepage
\begin{center}
\textbf{Goal:} connect surprise, uncertainty, and model mismatch
\end{center}
\end{frame}

\begin{frame}[c]{Surprise and information}
A rare event is more surprising than a common one.

\vspace{0.5em}
We want a function $I(x)$ with these properties:
\begin{enumerate}
    \item rarity gives larger information
    \item certainty gives zero information
    \item independent events add information
    \item small probability changes cause small information changes
\end{enumerate}

\vspace{0.5em}
These requirements force the logarithm:
\[
I(x) = -\log_2 p(x)
\]

\vspace{0.5em}
A likely event gives little surprise; an unlikely one gives a lot.
\end{frame}

\begin{frame}[c]{Self-information}
The self-information of an outcome is:
\[
I(x) = -\log_2 p(x)
\]

\vspace{0.5em}
Examples:
\begin{itemize}
    \item $p(x)=1/2 \Rightarrow I(x)=1$ bit
    \item $p(x)=1/4 \Rightarrow I(x)=2$ bits
    \item $p(x)=1 \Rightarrow I(x)=0$
\end{itemize}

\vspace{0.5em}
This gives a clean measure of surprise in bits.
\end{frame}

\begin{frame}[c]{Entropy: average surprise}
Entropy is the expected surprise under a distribution $p$:
\[
H(p) = -\sum_{x \in \mathcal{X}} p(x)\log_2 p(x)
\]

\vspace{0.5em}
\textbf{Interpretation:}
\begin{itemize}
    \item deterministic distribution $\Rightarrow H=0$
    \item more uncertainty $\Rightarrow$ larger entropy
    \item uniform distribution maximizes entropy
\end{itemize}
\end{frame}

\begin{frame}[c]{Maximum entropy}
For a finite set of $n$ outcomes, the uniform distribution maximizes entropy:
\[
H = -\sum_{i=1}^n \frac{1}{n}\log_2\left(\frac{1}{n}\right) = \log_2(n)
\]

\vspace{0.5em}
This is the maximum entropy principle: if we know only the set of possible outcomes, a uniform distribution introduces the fewest extra assumptions.
\end{frame}

\begin{frame}[c]{Cross-entropy}
Cross-entropy compares the real distribution $p$ with a model distribution $q$:
\[
H(p,q)= -\sum_{x\in\mathcal{X}} p(x)\log_2 q(x)
\]

\vspace{0.5em}
\textbf{Interpretation:}
\begin{itemize}
    \item if $q=p$, then $H(p,q)=H(p)$
    \item if the model is worse, cross-entropy is larger
\end{itemize}

\vspace{0.5em}
This is the training objective for language models: minimize surprise on real text.
\end{frame}

\begin{frame}[c]{KL divergence}
KL divergence tells us how much extra surprise comes from using $q$ instead of $p$:
\[
D_{KL}(p\|q)=\sum_x p(x)\log_2\frac{p(x)}{q(x)}
\]

\vspace{0.5em}
A useful identity is:
\[
D_{KL}(p\|q)=H(p,q)-H(p)
\]

\vspace{0.5em}
KL divergence is always nonnegative and is zero only when $p=q$.
\end{frame}

\begin{frame}[c]{Why direction matters}
KL divergence is not symmetric:
\[
D_{KL}(p\|q) \ne D_{KL}(q\|p)
\]

\vspace{0.5em}
This matters because each direction penalizes different mistakes:
\begin{itemize}
    \item $D_{KL}(p\|q)$ emphasizes missed real modes in $p$
    \item $D_{KL}(q\|p)$ emphasizes mass placed where $p$ has none
\end{itemize}

\vspace{0.5em}
In generative AI evaluation, we often choose the direction that matches the failure mode we care about.
\end{frame}

\begin{frame}[c]{Generative AI connection}
\textbf{Training:} language models minimize cross-entropy on text data.

\vspace{0.5em}
For one-hot targets, the loss is just:
\[
-\log q(\text{correct token})
\]

\vspace{0.5em}
\textbf{Evaluation:} lower cross-entropy means better predictive fit.

\vspace{0.5em}
\textbf{Fine-tuning:} KL penalties keep models from drifting too far from the base model.
\end{frame}

\begin{frame}[c]{Takeaway}
\textbf{Self-information} measures surprise of one event.

\textbf{Entropy} measures average surprise of a whole distribution.

\textbf{Cross-entropy} measures how surprised a model is by real data.

\textbf{KL divergence} measures the mismatch between two distributions.

\vspace{0.5em}
These ideas are at the heart of model training, evaluation, and uncertainty in generative AI.
\end{frame}

\end{document}
